% =========================================================
%  SECCIÓN III — METODOLOGÍA
% =========================================================
\section{Metodolog\'ia}

\subsection{Descripci\'on del Kernel}

El kernel num\'erico implementado calcula la magnitud del
gradiente del operador Prewitt (ec.~\ref{eq:magnitud}) sobre
un arreglo unidimensional de $N$ elementos de tipo
\texttt{float} (32~bits).
Las im\'agenes de entrada tienen un tama\~no de $256 \times 256$
p\'ixeles en formato PGM (\textit{Portable Gray Map}) binario
(P5), que es el formato de imagen en escala de grises sin
compresi\'on m\'as sencillo disponible en el est\'andar Netpbm.

Para obtener mediciones estables de rendimiento con Intel Advisor
se requiere que el n\'umero de operaciones sea suficientemente
grande~\cite{intel_advisor}.
Por ello, el kernel se aplica sobre $N_{\text{tot}}$ elementos:

\begin{equation}
  N_{\text{tot}} = 256 \times 256 \times 800 = 52{,}428{,}800
  \label{eq:ntot}
\end{equation}

Es decir, el kernel se ejecuta como si se procesaran 800 r\'eplicas
de la imagen original de forma contigua en memoria, garantizando
un n\'umero de iteraciones suficiente para que el perfilador
de Intel Advisor pueda medir con precisi\'on los trip counts,
la intensidad aritm\'etica y el modelo Roofline.

Para asegurar la correcta generaci\'on de instrucciones SIMD,
los tres arreglos ($G_x$, $G_y$ y $\textit{Mag}$) se alojan con
alineaci\'on a 32~Bytes mediante \texttt{\_mm\_malloc}:

\begin{lstlisting}[language=C++,caption={Alocaci\'on alineada de arreglos.}]
float* Gx  = (float*)_mm_malloc(N * sizeof(float), 32);
float* Gy  = (float*)_mm_malloc(N * sizeof(float), 32);
float* Mag = (float*)_mm_malloc(N * sizeof(float), 32);
\end{lstlisting}

El calificador \texttt{\_\_restrict} se usa en todas las versiones
para indicar al compilador que los punteros no se solapan en
memoria (\textit{no aliasing}), condici\'on necesaria para que
la vectorizaci\'on autom\'atica sea segura.

\subsection{Plataforma Experimental}

\textbf{Hardware}: Procesador Intel de 14 n\'ucleos f\'isicos
(gen.~Core de \'ultima generaci\'on), con soporte para extensiones
AVX2 y SSE4.2.
La arquitectura reportada por Intel Advisor incluye:
techo vectorial SP FMA Peak de $148.8$~GFLOPS (AVX-512 FMA),
techo de adici\'on vectorial SP de $74.55$~GFLOPS (AVX2/AVX-512),
techo escalar SP de $9.28$~GFLOPS,
ancho de banda L1 $\approx 413$~GB/s,
ancho de banda L2 $\approx 190$~GB/s,
ancho de banda L3 $\approx 65$--$73$~GB/s,
y ancho de banda DRAM de $24$--$27$~GB/s.

\textbf{Software}: Compilador Intel \texttt{icpx} (Intel oneAPI
DPC++/C++ Compiler 2025.3.2) sobre Linux (Fedora 43, kernel
6.18.12, 64~bits).
Sistema de profiling: Intel Advisor versi\'on incluida en oneAPI
2025.3.2.

Las opciones de compilaci\'on para cada versi\'on son:

\begin{center}
\footnotesize
\begin{tabular}{@{}ll@{}}
\toprule
\textbf{Versi\'on} & \textbf{Flags de compilaci\'on} \\
\midrule
Escalar    & \texttt{-g -O2 -no-vec -std=c++17 -fno-inline} \\
Autom\'atica & \texttt{-g -O3 -march=native -std=c++17 -fno-inline} \\
Guiada     & \texttt{-g -O3 -march=native -qopenmp -std=c++17 -fno-inline} \\
Expl\'icita & \texttt{-g -O3 -march=native -std=c++17 -fno-inline} \\
\bottomrule
\end{tabular}
\end{center}

En todas las versiones se a\~nade \texttt{-qopt-report=max}
\texttt{-qopt-report-phase=vec} para generar el reporte de
optimizaci\'on del compilador (\texttt{.optrpt}), que
documenta si cada bucle fue vectorizado y con qu\'e ancho
vectorial.

\subsection{Versiones Implementadas}

\textbf{Versi\'on 1 — Escalar (l\'inea base)}: El kernel
no utiliza ninguna instrucci\'on SIMD.
El flag \texttt{-no-vec} proh\'ibe la vectorizaci\'on
autom\'atica al compilador:

\begin{lstlisting}[language=C++,caption={Kernel escalar (versi\'on de referencia).}]
void kernel_escalar(
    const float* __restrict Gx,
    const float* __restrict Gy,
    float* __restrict Mag, int N)
{
    for (int i = 0; i < N; i++)
        Mag[i] = sqrtf(Gx[i]*Gx[i] + Gy[i]*Gy[i]);
}
\end{lstlisting}

\textbf{Versi\'on 2 — Autom\'atica}: C\'odigo id\'entico al
escalar, pero compilado con \texttt{-O3 -march=native}.
El compilador Intel \texttt{icpx} detecta que el bucle es
vectorizable (punteros con \texttt{\_\_restrict}, sin
dependencias entre iteraciones) y genera instrucciones AVX2:

\begin{lstlisting}[language=C++,caption={Kernel con vectorizaci\'on autom\'atica.}]
void kernel_auto(
    const float* __restrict Gx,
    const float* __restrict Gy,
    float* __restrict Mag, int N)
{
    for (int i = 0; i < N; i++)
        Mag[i] = sqrtf(Gx[i]*Gx[i] + Gy[i]*Gy[i]);
}
\end{lstlisting}

\textbf{Versi\'on 3 — Guiada}: Se a\~nade la directiva
\texttt{\#pragma omp simd simdlen(8)} para indicar expl\'icitamente
al compilador que vectorice con un ancho de 8~elementos
(256~bits, registros YMM):

\begin{lstlisting}[language=C++,caption={Kernel con vectorizaci\'on guiada.}]
void kernel_guiado(
    const float* __restrict Gx,
    const float* __restrict Gy,
    float* __restrict Mag, int N)
{
    #pragma omp simd simdlen(8)
    for (int i = 0; i < N; i++)
        Mag[i] = sqrtf(Gx[i]*Gx[i] + Gy[i]*Gy[i]);
}
\end{lstlisting}

\textbf{Versi\'on 4 — Expl\'icita}: El programador escribe
directamente las instrucciones AVX2 mediante intr\'insecos
\texttt{\_mm256\_*}.
El bucle principal procesa 8~floats por iteraci\'on; el bucle
residual maneja los elementos sobrantes cuando $N$ no es
m\'ultiplo de~8:

\begin{lstlisting}[language=C++,caption={Kernel con intr\'insecos AVX2 expl\'icitos.}]
void kernel_explicito(
    const float* __restrict Gx,
    const float* __restrict Gy,
    float* __restrict Mag, int N)
{
    int i = 0;
    for (; i <= N - 8; i += 8) {
        __m256 vx  = _mm256_loadu_ps(&Gx[i]);
        __m256 vy  = _mm256_loadu_ps(&Gy[i]);
        __m256 vx2 = _mm256_mul_ps(vx, vx);
        __m256 vy2 = _mm256_mul_ps(vy, vy);
        __m256 sum = _mm256_add_ps(vx2, vy2);
        __m256 res = _mm256_sqrt_ps(sum);
        _mm256_storeu_ps(&Mag[i], res);
    }
    for (; i < N; i++)  // residual escalar
        Mag[i] = sqrtf(Gx[i]*Gx[i] + Gy[i]*Gy[i]);
}
\end{lstlisting}

\subsection{Procedimiento de An\'alisis con Intel Advisor}

Para cada versi\'on se ejecutaron dos fases de an\'alisis de
Intel Advisor desde l\'inea de comandos:

\begin{enumerate}
  \item \textbf{Survey + mezcla de instrucciones est\'aticas}:
\begin{lstlisting}[language=bash,caption={Fase Survey de Intel Advisor.},numbers=none]
advisor --collect=survey
        --static-instruction-mix
        --project-dir=./adv_X
        -- ./ejecutable
\end{lstlisting}

  \item \textbf{Tripcounts + FLOP}:
\begin{lstlisting}[language=bash,caption={Fase Tripcounts de Intel Advisor.},numbers=none]
advisor --collect=tripcounts
        --flop
        --project-dir=./adv_X
        -- ./ejecutable
\end{lstlisting}
\end{enumerate}

Posteriormente se abri\'o la interfaz gr\'afica de Intel Advisor
(\texttt{advisor-gui}) con los proyectos generados para obtener
las gr\'aficas Roofline, los Program Metrics y el Survey de
bucles, de los que se extrajeron capturas de pantalla para su
an\'alisis.
